%% ==============================
% Part is used only in PhD thesis
\part{The Evaluation}
\chapter{\iflanguage{ngerman}{Diskussion}{Discussion}}
\label{sec:discussion}
%% ==============================

\section{Modellasymmetrie als durchgehender Erklärungsrahmen}
\label{sec:discussion-asymmetrie}

\begin{itemize}
  \item Ausgangspunkt: HCS objektgebundene, schlüsselbasierte Schreibkontrolle je Topic; Indy netzwerkweite, rollenbasierte Schreibkontrolle (Abschnitt~\ref{sec:impl-befunde}, Befund B4).
  \item Diese Asymmetrie wurde im Entwurf selbst antizipiert, nicht erst in der Auswertung entdeckt (Abschnitt~\ref{subsec:e6-vertrauenshierarchie}: „...zugleich entsteht die Asymmetrie gegenüber Indy, das den unberechtigten Schreibvorgang bereits abweist, die das Szenario S5 untersucht").
  \item Durchschlag in allen fünf Szenarien:
  \begin{itemize}
    \item S3: Zustandsableitung je Lesezugriff (HCS, E5) gegen natives Zustandsobjekt (Indy) - Ursache der (im getesteten Bereich nicht messbaren) Lesezeit-Bestandsgrößen-Abhängigkeit.
    \item S4: Für HCS liegt ein reproduzierter, bis auf Quellcode-Ebene erklärter Befund vor (fehlende Node-Fallback-Logik, sofortiges Scheitern). Für Indy zeigen mehrere unabhängige Nachprüfungsläufe eine reproduzierbare protokollinterne Selbstheilung nach Primary-Ausfall (View-Change, ca.\ eine Minute), belegt unter anderem durch direkte Client-Log-Analyse. Eine ältere Messserie unter identischer Konfiguration weicht davon vollständig ab (keine Erholung über 170\,s) und bleibt trotz gezielter Prüfung mehrerer Erklärungsansätze (Client-Verhalten, Primary-Status, Code-Änderungen, Beobachtungsfenster-Länge) ungeklärt -- ausgewiesen als Grenze der Aussagekraft dieser Messreihe.
  \end{itemize}
  \item Einordnung: kein allgemeines Überlegenheitsurteil, sondern eine strukturelle Weichenstellung, die die meisten beobachteten Differenzen erklärt.
\end{itemize}

\section{Wiederkehrendes Muster: Klientbibliothek begrenzt, was die Technologie nicht begrenzt}
\label{sec:discussion-bibliotheksgrenzen}

\begin{itemize}
  \item Vier unabhängige Instanzen desselben Grundmusters, alle in derselben Bibliotheksfamilie (\texttt{hiero\_did\_sdk\_python}/\texttt{hiero\_sdk\_python}):
  \begin{enumerate}
    \item Agentenpfad-Boden der Leseabonnement-Schnittstelle (S1, kausal bewiesen über Zeitfenster-Variation, Befund B5).
    \item Adressierung stets desselben einzelnen Knotens unabhängig von der Konfiguration (S4, Befund B5).
    \item Keine Wiederholung bei transienter \texttt{BUSY}-Ablehnung im Schreibpfad, obwohl Lesepfad und Quittungsabfragen dieselbe Bibliotheksversion wiederholen (S2, Befund B6).
    \item (schwächer belegt) zweigipfliges Zeitmuster im Indy-CredDef-Schreibpfad (S1), last-abhängiger Anteil nicht bis zum Quellcode zurückverfolgt.
  \end{enumerate}
  \item Gemeinsamkeit: systematisches Reifegradgefälle zwischen Lese- und Schreibpfad in dieser SDK-Generation, kein Einzelfall.
  \item Zeitlich begrenzt: mindestens drei der vier Fälle sind in neueren Versionen derselben Bibliothek behoben; keiner wurde vor dieser Arbeit gemeldet.
  \item Konsequenz für die Gesamtinterpretation: die gemessenen HCS-Nachteile sind überwiegend dem zum Messzeitpunkt amtierenden Software-Stack zuzuschreiben, nicht dem Hedera-Konsensprotokoll selbst. Eine Aussage über die Technologie „Hashgraph" und eine Aussage über „diese Implementierung, Stand 2026" müssen getrennt gelesen werden.
\end{itemize}

\section{Einordnung je nicht-funktionaler Anforderung}
\label{sec:discussion-nfa}

\begin{itemize}
  \item \textbf{NFA1/NFA2 (S1, S2):} % TODO kurzer Absatz, Zahlen aus Kapitel 6 aufgreifen
  \item \textbf{NFA4 (S2):} Anforderung verlangt Sättigungspunkt \emph{und} begrenzende Komponente (Abschnitt~\ref{subsec:eval-messgroessen}) - für HCS beides identifiziert (Fehlerrate ab 10 Nutzern, Ursache im Schreibpfad der Klientbibliothek), für Indy ebenso (Durchsatz-Plateau bei 55 Nutzern).
  \item \textbf{NFA5 (S3):} Nullergebnis als obere Schranke, nicht als Widerlegung des E5-Arguments; Konzept selbst sagt vorher, dass der Effekt erst bei größeren Beständen sichtbar wird.
  \item \textbf{NFA6 (S4):} wichtige Differenzierung - gemessen wurde die Unfähigkeit des amtierenden Software-Stacks, HCS' aBFT-Eigenschaften zugänglich zu machen, nicht die Abwesenheit dieser Eigenschaften selbst. Die Frage „ist Hedera-Hashgraph fehlertolerant" bleibt technisch unbeantwortet; beantwortet ist „ist dieser Stack, Stand 2026, fehlertolerant nutzbar" - mit nein. Methodische Konsequenz: durch die fest verdrahtete Einzelknoten-Adressierung ergibt das Szenario auf HCS strukturell nur zwei mögliche Ausgänge, kein Kontinuum - entweder der tatsächlich angesprochene Knoten fällt aus (Erfolgsrate nahe null) oder ein anderer (kein messbarer Effekt); ein vierknotiges Netzwerk liefert unter diesem Software-Stand keinen realisierten Mehrwert gegenüber einem einzelnen Knoten, obwohl die Redundanz betrieben und bezahlt wird.
  \item \textbf{NFA7 (S5):} zweigeteilt einzulösen (Abschnitt~\ref{subsec:eval-validitaet}). Netzwerkseitige Zugriffskontrolle: auf beiden Varianten vollständig gemessen. Anwendungsseitige Stufe: durch die Nichtrealisierung von E6 (Abschnitt~\ref{sec:impl-befunde}, Befund B4) nicht auf HCS testbar; als Gültigkeitsgrenze zu benennen, nicht zu verschweigen. Widerrufs-Mindestanforderung: unabhängig davon auf beiden Varianten vollständig erfüllt.
\end{itemize}

\section{Grenzen der Aussagekraft}
\label{sec:discussion-grenzen}

\begin{itemize}
  \item Grundlast bewusst synthetisch/definiert: Taktrate als gesetzter Versuchsparameter, nicht aus beobachtetem Nutzungsverhalten abgeleitet (Abschnitt~\ref{subsec:eval-szenarien}); ein empirisch fundiertes Lastprofil für die Referenzarchitektur liegt bislang nicht vor.
  \item E5-Leseweg bewusst nicht optimiert; berichtete HCS-Lesewerte sind obere Schranke des Aufwands, keine Leistungsgrenze der Technologie (Zitat Konzept, Abschnitt~\ref{sec:konzipierung-zielarchitektur}).
  \item S3-Bestandsgrößen (bis 300 Objekte) möglicherweise zu klein, um einen tatsächlich vorhandenen Effekt sichtbar zu machen.
  \item S5-Indy-Widerrufsmessung und S5-Indy-Endorser-Messung beide ohne Vier-Knoten-Gegenprobe (im Unterschied zur HCS-Widerrufsmessung, die durch unabhängige Mirror-Node-Historie bestätigt ist) - Einzelmessungen, nicht durch unabhängige Knotenprotokolle bestätigt.
  \item Anwendungsseitige E6-Stufe nicht realisiert (siehe oben) - S5 auf HCS strukturell unvollständig gegenüber der kanonischen Definition.
\end{itemize}

\section{Gesamtbild}
\label{sec:discussion-gesamtbild}

\begin{itemize}
  \item Funktionale Substitution gelungen: alle fünf Ledger-Funktionen (DID-Registrierung, Schema, CredDef, Widerruf, Zugriffskontrolle) mit HCS-Mechanismen nachgebildet und in S1-S5 erfolgreich durchgeführt; die in Kapitel~\ref{cha:konzipierung} entworfene Zielarchitektur ist tragfähig.
  \item Die gemessenen Nachteile sind überwiegend Eigenschaften der eingesetzten Klientbibliotheken, nicht des Hashgraph-Konsens selbst (Abschnitt~\ref{sec:discussion-bibliotheksgrenzen}); dieses Muster trägt drei der vier Hauptbefunde (S1-Boden, S2-Kollaps, S4-Nichtausfall) und rechtfertigt eine getrennte Lesart von "Technologie Hashgraph" und "amtierender Software-Stand 2026".
  \item Der architektonische Preis der in E5 gewählten Zustandsableitung beim Lesen ist real, aber im getesteten Rahmen nicht dramatisch: S3 findet bis 300 Füllobjekten keinen messbaren Bestandsgrößeneffekt - eine obere, keine endgültige Schranke (Abschnitt~\ref{sec:discussion-nfa}, NFA5).
  \item Sicherheitsseitig ist Indys Rollenmodell reicher und beobachtbarer, HCS' Modell deshalb aber nicht schwächer, sondern andersartig: objektgebundene Schlüsselkontrolle statt Netzwerk-Rollenhierarchie (Befund B4); die fehlende anwendungsseitige Delegationsstufe (E6) ist ein unimplementierter, kein protokollseitiger Nachteil.
  \item Evidenzstärke variiert bewusst über die Befunde und wird durchgehend so ausgewiesen: kausal bewiesen (S1-HCS-Zeitfenster, S2-\texttt{BUSY}-Pfad) versus beobachtet, nicht mechanistisch belegt (S1-Indy-Bimodalität) versus obere Schranke statt Nullbefund (S3); diese Abstufung wird nicht zu einer einheitlichen Sicherheitsstufe verflacht.
\end{itemize}

\section{Ausblick}
\label{sec:discussion-ausblick}

\begin{itemize}
  \item Größere S3-Bestandsserie zur Absicherung des Nullergebnisses.
  \item S2/S4 mit retry-fähiger Klientbibliothek als separat gelabelter Interventions-Arm - Kosten/Nutzen-Abwägung dokumentiert, im Rahmen dieser Arbeit nicht verfolgt.
  \item Anwendungsseitige E6-Stufe als eigenständige Implementierungsarbeit für eine Folgearbeit.
  \item % TODO weitere Punkte nach Fertigstellung von Kapitel 6 ergaenzen
\end{itemize}